\chapter{Protokol MQTT}
\label{ch:06-protokol-mqtt}

\section{Kerangka Topik dan Penamaan}

Seluruh pesan bergerak pada topik dengan tiga segmen: prefix proyek, identitas
perangkat, dan jenis pesan. Prefix bawaan adalah \code{\TopicPrefix}, identitas
perangkat diturunkan dari MAC chip sehingga berbentuk
\code{<\DeviceIdPrefix>-<tiga byte terakhir>} seperti \code{\DeviceExample}, dan
jenis pesan adalah salah satu dari \code{status}, \code{telemetry}, \code{ack},
\code{state}, atau \code{command}. Contoh topik lengkap untuk perangkat contoh
adalah \topic{telemetry}. Baris perangkat awal di basis data memakai identitas
\code{\SeedDevice}, sehingga topiknya berbentuk \topicof{\SeedDevice/ack}.

Tiga segmen itu bukan sekadar konvensi penamaan. Worker platform memecah topik
menjadi tepat tiga bagian dan membuang pesan yang jumlah segmennya berbeda atau
yang prefix-nya tidak cocok, sehingga topik dari proyek lain pada broker yang sama
tidak pernah masuk ke basis data. Perangkat melakukan penyaringan sebaliknya:
hanya pesan pada topik perintah miliknya sendiri yang diproses, dan seluruh pesan
lain diabaikan tanpa jawaban.

\section{Tabel Topik}

\tabref{tab:06-topik} adalah rujukan tunggal untuk semua topik yang dipakai
sistem. Perhatikan bahwa hanya topik status yang memakai pesan retained; sisanya
tidak, karena telemetri, ack, dan event bersifat peristiwa yang tidak boleh
diputar ulang kepada subscriber baru.

\begin{table}[H]
\centering
\caption{Topik MQTT, arah, tanda retained, QoS, dan isi payload.}
\label{tab:06-topik}
\small
\begin{tabular}{@{}p{1.6cm}p{2.7cm}p{1.2cm}p{1.5cm}p{6.9cm}@{}}
\toprule
\textbf{Kind} & \textbf{Arah} & \textbf{Retain} & \textbf{QoS} & \textbf{Payload} \\
\midrule
\code{status} & perangkat ke platform & ya & 0 & String \code{online} saat
connect dan \code{offline} sebagai LWT. Bukan JSON. \\
\code{telemetry} & perangkat ke platform & tidak & 0 & JSON berisi identitas,
state, jarak, persentase, dan waktu. \\
\code{ack} & perangkat ke platform & tidak & 0 & JSON berisi
\code{command_id}, hasil, pesan, dan waktu. \\
\code{state} & perangkat ke platform & tidak & 0 & JSON event
\code{scheduled_feed} yang dipicu perangkat sendiri. \\
\code{command} & platform ke perangkat & tidak & 1 & JSON berisi
\code{command_id}, tipe, payload, dan waktu. \\
\bottomrule
\end{tabular}
\end{table}

\src{firmware/esp32_feeder/net.cpp, dashboard/mqtt.py}

Asimetri QoS perlu dicatat karena sering disalahpahami. Perangkat mempublikasikan
dengan QoS~0 tetapi berlangganan perintah dengan QoS~1; platform mempublikasikan
perintah dengan QoS~1 tetapi berlangganan dengan QoS~0. Artinya jaminan
pengiriman hanya diberikan pada arah perintah, yaitu arah yang tidak boleh hilang.
Pada arah telemetri, kehilangan satu sampel tidak merusak apa pun karena sampel
berikutnya datang \TelemetryMs{} ms kemudian, dan keandalan yang sebenarnya
diperiksa pada lapisan aplikasi lewat ack, bukan pada lapisan transport.

\section{Contoh Payload}

Payload telemetri memuat jarak mentah dan persentase hasil hitung firmware. Field
yang nilainya tidak diketahui dihilangkan, bukan diisi nol, sehingga platform dapat
membedakan ``jarak tidak terbaca'' dari ``hopper kosong''.

\begin{lstlisting}[style=report]
{"device":"feeder-64008c","state":"idle","distance_cm":10.0,
 "feed_level_pct":50,"ts":"2026-11-19T13:00:00"}
\end{lstlisting}

Nilai \code{state} hanya boleh salah satu dari \code{idle}, \code{feeding}, atau
\code{error}. Nilai \code{error} dipakai ketika pembacaan jarak tidak valid,
sehingga dashboard dapat menampilkan keadaan yang jujur alih-alih angka yang
kebetulan tersisa dari pengukuran lama.

Event jadwal dipublikasikan setelah perangkat mengeksekusi jadwalnya sendiri.
Jarak pada event diukur sesudah pemberian pakan, dan justru itu gunanya: platform
memakai baris telemetri terakhir sebagai level sebelum, dan jarak pada event
sebagai level sesudah.

\begin{lstlisting}[style=report]
{"event":"scheduled_feed","device":"feeder-64008c","portions":2,
 "feed_time":"07:00:00","distance_cm":14.8,"ts":"2026-11-19T07:00:01"}
\end{lstlisting}

Ack adalah pesan yang mengikat perintah dengan hasilnya. Field \code{command_id}
harus sama persis dengan yang diterima, karena nilai itulah satu-satunya kunci
korelasi yang dimiliki platform.

\begin{lstlisting}[style=report]
{"command_id":"cmd-ed239f2b-f305-46e6-8998-4fa7a4492986",
 "result":"ok","message":"feeding 2 portions","ts":"2026-11-19T13:50:19"}
\end{lstlisting}

Perintah yang dikirim platform membawa tipe dan payload bertingkat, sehingga
penambahan field pada satu tipe tidak mengubah bentuk pesan tipe lain.

\begin{lstlisting}[style=report]
{"command_id":"cmd-ed239f2b-f305-46e6-8998-4fa7a4492986",
 "type":"manual_feed","payload":{"portions":2},
 "timestamp":"2026-11-19T13:50:18Z"}
\end{lstlisting}

\section{Perintah dan Ack}

Ada enam tipe perintah, dan setiap tipe memiliki urutan pemeriksaan yang tetap.
Urutan itu bukan detail implementasi: ia menentukan kode hasil mana yang diterima
pengguna ketika dua masalah terjadi sekaligus. Sebagai contoh, permintaan pakan
saat gate masih bergerak selalu dijawab \code{error_busy}, bukan
\code{error_low_level}, karena pemeriksaan kesibukan dilakukan lebih dahulu.
\tabref{tab:06-command} memuat keenam tipe beserta payload, urutan penjagaan, dan
hasil suksesnya.

\begin{table}[H]
\centering
\caption{Perintah, payload, urutan penjagaan, dan hasil yang diharapkan.}
\label{tab:06-command}
\small
\begin{tabular}{@{}p{2.7cm}p{2.8cm}p{4.9cm}p{3.9cm}@{}}
\toprule
\textbf{Tipe} & \textbf{Payload} & \textbf{Penjagaan (berurutan)} & \textbf{Hasil sukses} \\
\midrule
\code{servo_test} & \code{angle} bernilai 0 atau 90 & Sudut selain 0 atau 90
menghasilkan \code{error_bad_payload}; gate sibuk menghasilkan
\code{error_busy}; servo tidak terpasang menghasilkan \code{error_servo}. &
\code{ok} dengan pesan sudut yang dituju. \\
\code{manual_feed} & \code{portions} 1 sampai 10 & Gate sibuk menghasilkan
\code{error_busy}; servo tidak terpasang menghasilkan \code{error_servo}; level
nol menghasilkan \code{error_low_level}. & \code{ok} setelah porsi terakhir
selesai. \\
\code{set_schedule} & Array jadwal aktif & Payload yang gagal di-parse
menghasilkan \code{error_bad_payload}. & \code{ok} dengan jumlah jadwal
tersimpan. \\
\code{set_threshold} & \code{threshold_pct} 0 sampai 100 & Nilai di luar rentang
menghasilkan \code{error_bad_payload}. & \code{ok} dengan ambang yang dipakai. \\
\code{set_rtc} & \code{datetime} berbentuk ISO tanpa zona waktu & Format yang
gagal di-parse menghasilkan \code{error_bad_payload}. & \code{ok}; bila RTC tidak
ada, pesan menyebutkan pemakaian jam internal. \\
\code{get_state} & Objek kosong & Tidak ada penjagaan khusus. & \code{ok} dengan
snapshot keadaan di dalam field \code{message}. \\
\bottomrule
\end{tabular}
\end{table}

\src{firmware/esp32_feeder/esp32_feeder.ino, firmware/esp32_feeder/README.md}

Empat aturan berikut berlaku untuk semua tipe tanpa kecuali.

\begin{enumerate}[leftmargin=*,itemsep=2pt]
  \item \textbf{Ack wajib.} Setiap perintah yang diterima harus dijawab, baik
        berhasil maupun gagal. Diam bukan jawaban yang sah: perangkat yang diam
        membuat pengguna menunggu tanpa kepastian selama \CmdTimeoutS{} detik.
  \item \textbf{\code{command_id} di-echo apa adanya.} Platform tidak memakai
        urutan kedatangan atau isi payload untuk mencocokkan jawaban; hanya
        \code{command_id} yang dipakai.
  \item \textbf{Batas waktu \CmdTimeoutS{} detik.} Jawaban yang datang setelah
        batas itu sudah tidak relevan karena perintahnya sudah ditandai
        \code{timeout} oleh sweeper.
  \item \textbf{Hanya topik perintah yang diproses.} Pesan pada topik lain
        diabaikan tanpa jawaban, termasuk pesan retained milik topik status yang
        mungkin ikut terkirim oleh broker tertanam.
\end{enumerate}

\section{Kode Hasil dan Artinya bagi Platform}

\tabref{tab:06-result} memetakan setiap kode hasil ke tindakan platform. Yang perlu
diperhatikan adalah bahwa semua kode selain \code{ok} diperlakukan sama oleh
platform, yaitu kegagalan yang disertai notifikasi; pemilahan lebih lanjut hanya
terjadi pada teks pesan yang disimpan untuk penelusuran.

\begin{table}[H]
\centering
\caption{Kode hasil pada ack dan tindakan platform.}
\label{tab:06-result}
\small
\begin{tabular}{@{}p{3.4cm}p{4.6cm}p{7.0cm}@{}}
\toprule
\textbf{Kode hasil} & \textbf{Arti} & \textbf{Tindakan platform} \\
\midrule
\code{ok} & Perintah dijalankan & Perintah dan riwayat pakan ditandai sukses;
level sesudah diambil dari telemetri terakhir. \\
\code{error_bad_payload} & Payload tidak sesuai kontrak & Perintah ditandai
gagal; pesan perangkat disimpan pada kolom hasil. \\
\code{error_busy} & Gate sedang menjalankan porsi lain & Ditandai gagal; pengguna
dapat mengirim ulang setelah pemberian pakan sebelumnya selesai. \\
\code{error_servo} & Servo tidak terpasang & Ditandai gagal dan memunculkan
notifikasi pemberian pakan gagal. \\
\code{error_low_level} & Hopper terbaca kosong & Ditandai gagal dan
memunculkan notifikasi pemberian pakan gagal. \\
\code{error_unknown_type} & Tipe perintah tidak dikenal perangkat & Ditandai
gagal; menunjukkan ketidaksesuaian versi antara platform dan firmware. \\
\bottomrule
\end{tabular}
\end{table}

Pesan yang tidak pernah sampai ke perangkat tidak menghasilkan kode apa pun.
Kasus itu seluruhnya ditangani oleh sweeper: perintah berubah menjadi
\code{timeout} dan riwayat pakan terkait menjadi gagal dengan pesan
\code{timeout}. Dengan pembagian ini, setiap perintah yang pernah dikirim platform
berakhir pada salah satu dari empat status, yaitu sukses, gagal, timeout, atau
masih berjalan, dan tidak ada perintah yang menggantung tanpa status.

\section{Batas Kontrak}

Kontrak ini tidak memiliki mekanisme idempotensi. Bila perangkat menerima
perintah yang sama dua kali, misalnya karena pengiriman ulang setelah koneksi
terputus, gate akan bergerak dua kali. Rancangan mengandalkan tiga hal untuk
membatasinya: QoS~1 hanya pada arah perintah, uuid per perintah yang dibuat baru
setiap kali pengguna menekan tombol, dan penjagaan kesibukan pada gate yang
menolak perintah kedua selama porsi pertama masih berjalan. Ketiganya mengurangi
kemungkinan, bukan menghilangkannya, dan keterbatasan ini dicatat pada
\chapref{ch:09-catatan-asumsi}.
